%! Direct, Inverse \& Joint Variation — Unit 4, Lesson 4.7 — Group Activity (Answer Key)
\documentclass[10pt]{article}
\usepackage{algebra2-article}
\usepackage{algebra2-key}   % pulls in -boxes; do NOT also load -boxes
\newcommand{\TallMath}[1]{$\displaystyle #1\rule[-1.4em]{0pt}{3.2em}$}

\begin{document}

\pageheader{Unit 4, Lesson 4.7}{Group Activity --- Answer Key}
\namepartnerperiod

\vspace{0.08in}

\begin{headlinebox}{forestbg}
\small\textbf{Constant Quotient, Constant Product.} One question decides every table:
\textbf{which row of arithmetic stays the same?} Quotients constant $\Rightarrow$ \emph{direct},
$y=kx$. Products constant $\Rightarrow$ \emph{inverse}, $y=\frac kx$. Neither constant
$\Rightarrow$ \emph{neither} --- a real answer, and a common one. Test \emph{every} row, not the
first one.
\end{headlinebox}

\vspace{0.06in}

\begin{tcolorbox}[colback=white, colframe=black!40, title=\textbf{Tier R --- Remediate},
    fonttitle=\bfseries, arc=1.5mm, left=3mm, right=3mm, top=2mm, bottom=2mm, breakable]
\textbf{Part 1 --- Run both tests.} Fill in both arithmetic rows, then circle your verdict and give
$k$ if there is one.
\begin{center}
\begin{minipage}{0.47\linewidth}\centering
\renewcommand{\arraystretch}{1.25}
\begin{tabular}{|l|c|c|c|c|}
\hline
\rowcolor{forestbg}
$x$ & $3$ & $5$ & $8$ & $10$ \\ \hline
$y$ & $12$ & $20$ & $32$ & $40$ \\ \hline
$y\div x$ & \ans{$4$} & \ans{$4$} & \ans{$4$} & \ans{$4$} \\ \hline
$x\cdot y$ & \ans{$36$} & \ans{$100$} & \ans{$256$} & \ans{$400$} \\ \hline
\end{tabular}\\[2pt]
{\small (a) \; \textcolor{keyred}{\textbf{direct}} \; / \; inverse \; / \; neither \quad
$k=$ \ans{$4$}}
\end{minipage}\hfill
\begin{minipage}{0.47\linewidth}\centering
\renewcommand{\arraystretch}{1.25}
\begin{tabular}{|l|c|c|c|c|}
\hline
\rowcolor{forestbg}
$x$ & $2$ & $4$ & $8$ & $16$ \\ \hline
$y$ & $24$ & $12$ & $6$ & $3$ \\ \hline
$y\div x$ & \ans{$12$} & \ans{$3$} & \ans{$0.75$} & \ans{$0.1875$} \\ \hline
$x\cdot y$ & \ans{$48$} & \ans{$48$} & \ans{$48$} & \ans{$48$} \\ \hline
\end{tabular}\\[2pt]
{\small (b) \; direct \; / \; \textcolor{keyred}{\textbf{inverse}} \; / \; neither \quad
$k=$ \ans{$48$}}
\end{minipage}
\end{center}
\vspace{4pt}
\begin{center}
\begin{minipage}{0.47\linewidth}\centering
\renewcommand{\arraystretch}{1.25}
\begin{tabular}{|l|c|c|c|c|}
\hline
\rowcolor{forestbg}
$x$ & $1$ & $2$ & $3$ & $4$ \\ \hline
$y$ & $20$ & $17$ & $14$ & $11$ \\ \hline
$y\div x$ & \ans{$20$} & \ans{$8.5$} & \ans{$4.\overline{6}$} & \ans{$2.75$} \\ \hline
$x\cdot y$ & \ans{$20$} & \ans{$34$} & \ans{$42$} & \ans{$44$} \\ \hline
\end{tabular}\\[2pt]
{\small (c) \; direct \; / \; inverse \; / \; \textcolor{keyred}{\textbf{neither}} \quad
$k=$ \ans{none}}
\end{minipage}\hfill
\begin{minipage}{0.47\linewidth}\centering
\renewcommand{\arraystretch}{1.25}
\begin{tabular}{|l|c|c|c|c|}
\hline
\rowcolor{forestbg}
$x$ & $2$ & $3$ & $5$ & $7$ \\ \hline
$y$ & $14$ & $21$ & $35$ & $49$ \\ \hline
$y\div x$ & \ans{$7$} & \ans{$7$} & \ans{$7$} & \ans{$7$} \\ \hline
$x\cdot y$ & \ans{$28$} & \ans{$63$} & \ans{$175$} & \ans{$343$} \\ \hline
\end{tabular}\\[2pt]
{\small (d) \; \textcolor{keyred}{\textbf{direct}} \; / \; inverse \; / \; neither \quad
$k=$ \ans{$7$}}
\end{minipage}
\end{center}
\vspace{2pt}
{\small In (c) the $y$-values go \emph{down}. That is not enough to make it inverse --- check the
product row and see.}

\vspace{5pt}
\textbf{Part 2 --- Find $k$, then answer.} Two steps every time.
\begin{enumerate}[label=(\alph*), leftmargin=2.1em, itemsep=5pt]
  \item $y$ varies \textbf{directly} as $x$, and $y=15$ when $x=3$. \; $k=$ \ans{$5$},
        equation $y=$ \ans{$5x$}. \; Find $y$ when $x=7$: \ans{$35$}
  \item $y$ varies \textbf{inversely} as $x$, and $y=9$ when $x=4$. \; $k=$ \ans{$36$},
        equation $y=$ \ans{$\dfrac{36}{x}$}. \; Find $y$ when $x=6$: \ans{$6$}
  \item $y$ varies \textbf{directly} as $x$, and $y=30$ when $x=12$. \; $k=$ \ans{$2.5$}. \;
        Find $x$ when $y=45$: \ans{$18$}
  \item $y$ varies \textbf{inversely} as $x$, and $y=5$ when $x=8$. \; $k=$ \ans{$40$}. \;
        Find $x$ when $y=2$: \ans{$20$}
\end{enumerate}

\vspace{5pt}
\textbf{Part 3 --- Twins.} In \emph{both} of these tables, $y$ goes up by exactly $5$ every time $x$
goes up by $2$. Exactly one of them is a direct variation.
\begin{center}
\begin{minipage}{0.47\linewidth}\centering
\textbf{Twin (a)}\\[2pt]
\renewcommand{\arraystretch}{1.25}
\begin{tabular}{|l|c|c|c|c|}
\hline
\rowcolor{forestbg}
$x$ & $2$ & $4$ & $6$ & $8$ \\ \hline
$y$ & $5$ & $10$ & $15$ & $20$ \\ \hline
$y\div x$ & \ans{$2.5$} & \ans{$2.5$} & \ans{$2.5$} & \ans{$2.5$} \\ \hline
\end{tabular}\\[3pt]
{\small Verdict \ans{direct}; \; equation \ans{$y=2.5x$}}
\end{minipage}\hfill
\begin{minipage}{0.47\linewidth}\centering
\textbf{Twin (b)}\\[2pt]
\renewcommand{\arraystretch}{1.25}
\begin{tabular}{|l|c|c|c|c|}
\hline
\rowcolor{forestbg}
$x$ & $2$ & $4$ & $6$ & $8$ \\ \hline
$y$ & $8$ & $13$ & $18$ & $23$ \\ \hline
$y\div x$ & \ans{$4$} & \ans{$3.25$} & \ans{$3$} & \ans{$2.875$} \\ \hline
\end{tabular}\\[3pt]
{\small Verdict \ans{neither}; \; equation \ans{$y=2.5x+3$}}
\end{minipage}
\end{center}
\vspace{-2pt}
Both tables climb at the same steady rate. What made the verdicts different? \ansline{Twin (b) starts
$3$ ahead: at $x=0$ it would give $y=3$, not $0$. A constant step makes a \emph{line}, but only a line
through the origin has a constant quotient}
\end{tcolorbox}

\begin{tcolorbox}[colback=white, colframe=black!40, title=\textbf{Tier A --- Approaching Proficiency},
    fonttitle=\bfseries, arc=1.5mm, left=3mm, right=3mm, top=2mm, bottom=2mm, breakable]
\textbf{Part 1 --- A spring.} Hang a weight on a spring and it stretches. Let $w$ be the weight in
pounds and $s$ the stretch in centimeters.
\begin{center}
\renewcommand{\arraystretch}{1.3}
\begin{tabular}{|l|c|c|c|c|}
\hline
\rowcolor{forestbg}
$w$ (lb) & $4$ & $6$ & $10$ & $15$ \\ \hline
$s$ (cm) & $3$ & $4.5$ & $7.5$ & $11.25$ \\ \hline
$s\div w$ & \ans{$0.75$} & \ans{$0.75$} & \ans{$0.75$} & \ans{$0.75$} \\ \hline
\end{tabular}
\end{center}
\begin{enumerate}[label=\arabic*., leftmargin=2.1em, itemsep=5pt]
  \item Verdict: \ans{direct} variation. \; $k=$ \ans{$0.75$}. \; Equation: $s=$ \ans{$0.75w$}
  \item What are the \emph{units} of $k$? \ans{cm per pound} \; Say what $k$ means in one sentence.
        \ans{each additional pound stretches the spring another $0.75$ cm}
  \item A $20$-pound weight stretches the spring \ans{$15$} cm. \; A stretch of $9$ cm means a
        weight of \ans{$12$} pounds.
  \item The graph of $s=kw$ passes through $(0,0)$. Explain why the \emph{situation} requires that.
        \ansline{with no weight hanging on it, the spring is not stretched at all --- zero weight must
        give zero stretch}
\end{enumerate}

\vspace{5pt}
\textbf{Part 2 --- Every rectangle with area $36$.} A rectangle has base $b$ and height $h$, and its
area is always $36$ square inches.
\begin{center}
\begin{minipage}{0.44\linewidth}\centering
\begin{tikzpicture}[scale=0.30]
  \draw[step=1, linegray!40, very thin] (0,0) grid (13,14);
  \draw[->, semithick, charcoal] (0,0)--(13.6,0) node[below right, font=\scriptsize]{$b$};
  \draw[->, semithick, charcoal] (0,0)--(0,14.6) node[left, font=\scriptsize]{$h$};
  \node[below, font=\tiny] at (3,0) {$3$};
  \node[below, font=\tiny] at (6,0) {$6$};
  \node[below, font=\tiny] at (9,0) {$9$};
  \node[below, font=\tiny] at (12,0) {$12$};
  \node[left, font=\tiny] at (0,3) {$3$};
  \node[left, font=\tiny] at (0,6) {$6$};
  \node[left, font=\tiny] at (0,9) {$9$};
  \node[left, font=\tiny] at (0,12) {$12$};
  \begin{scope}
    \clip (0,0) rectangle (13,14);
    \draw[very thick, forest, domain=2.6:13, samples=170, smooth] plot (\x,{36/(\x)});
  \end{scope}
  \fill[charcoal] (3,12) circle (3.4pt);
  \fill[charcoal] (6,6) circle (3.4pt);
  \fill[charcoal] (12,3) circle (3.4pt);
\end{tikzpicture}
\end{minipage}\hfill
\begin{minipage}{0.52\linewidth}
\small
\begin{enumerate}[label=(\alph*), leftmargin=1.9em, itemsep=4pt]
  \item $b$ and $h$ vary \ans{inversely}, because $b\cdot h$ is always \ans{$36$}. Equation:
        $h=$ \ans{$\dfrac{36}{b}$}
  \item Read the graph: when $b=6$, $h=$ \ans{$6$}. What is special about that rectangle?
        \ans{it is a square}
  \item Solve for the base when $h=8$: $b=$ \ans{$4.5$}. Is that point on the drawn curve?
        \ans{yes}
  \item The curve gets very close to the $b$-axis on the right but never touches it. Say what that
        means about rectangles. \ans{a rectangle of area $36$ can be made as short as you like by
        making it long enough, but its height can never actually reach $0$}
  \item The equation is also happy with $b=-4$. Why does the situation say no?
        \ans{a base is a length, so it has to be positive}
\end{enumerate}
\end{minipage}
\end{center}

\vspace{5pt}
\textbf{Part 3 --- Error analysis.} A student is handed this table and writes what follows.
\begin{center}
\renewcommand{\arraystretch}{1.25}
\begin{tabular}{|l|c|c|c|c|}
\hline
\rowcolor{forestbg}
$x$ & $2$ & $4$ & $6$ & $8$ \\ \hline
$y$ & $6$ & $10$ & $14$ & $18$ \\ \hline
\end{tabular}
\end{center}
\vspace{2pt}
\begin{center}
\fbox{\parbox{0.80\linewidth}{\centering \vspace{2pt}
``$x=2$ goes with $y=6$, so $k=\frac{6}{2}=3$. \\
The equation is $y=3x$. \\
It's a direct variation with $k=3$.''\vspace{2pt}}}
\end{center}
\begin{enumerate}[label=(\alph*), leftmargin=2.1em, itemsep=5pt]
  \item Is the student's arithmetic on the \emph{first} pair correct? \ans{yes}
  \item Test their equation on the rest of the table: $y=3x$ predicts \ans{$12$} at $x=4$, but
        the table says \ans{$10$}.
  \item Finish both rows properly. \; $y\div x$: \ans{$3,\;2.5,\;2.\overline{3},\;2.25$} \\
        $x\cdot y$: \ans{$12,\;40,\;84,\;144$} \; So the correct verdict is \ans{neither}.
  \item The table is $y=2x+2$. Write the one-sentence rule the student needed. \ansline{One pair of
        values can never decide the type --- every pair has both a quotient and a product.}
        \ansline{The constant has to hold for \emph{every} row in the table, so you check them all
        before you name it.}
\end{enumerate}
\end{tcolorbox}

\begin{tcolorbox}[colback=white, colframe=black!40, title=\textbf{Tier E --- Extension},
    fonttitle=\bfseries, arc=1.5mm, left=3mm, right=3mm, top=2mm, bottom=2mm, breakable]
\textbf{Part 1 --- Two variables at once} \emph{(joint and combined variation --- enrichment)}.
\begin{enumerate}[label=(\alph*), leftmargin=2.1em, itemsep=5pt]
  \item $y$ varies jointly as $x$ and $z$, and $y=48$ when $x=4$ and $z=2$. \; $k=$ \ans{$6$},
        equation $y=$ \ans{$6xz$}. \; Find $y$ when $x=5$, $z=3$: \ans{$90$}
  \item The volume $V$ of a cylinder varies jointly as its height $h$ and the \emph{square} of its
        radius $r$. When $r=5$ and $h=4$, $V=100\pi$. \\
        Then $k=$ \ans{$\pi$} and the equation is $V=$ \ans{$\pi r^2h$}. Do you recognize it?
        \ans{it is the volume formula for a cylinder}
  \item $y$ varies directly as $x$ and inversely as $z$, and $y=14$ when $x=7$ and $z=3$. \;
        $k=$ \ans{$6$}, equation $y=$ \ans{$\dfrac{6x}{z}$}. \; Find $y$ when $x=10$, $z=4$:
        \ans{$15$}
\end{enumerate}

\vspace{5pt}
\textbf{Part 2 --- No numbers required.} Fill the table using the equations alone.
\begin{center}
\renewcommand{\arraystretch}{1.5}
\begin{tabularx}{\linewidth}{@{}l l X@{}}
\hline
\textbf{Relationship} & \textbf{What you do to $x$ (and $z$)} & \textbf{What happens to $y$} \\
\hline
$y=kx$ & triple $x$ & \ans{$y$ is tripled} \\
$y=kx$ & cut $x$ in half & \ans{$y$ is cut in half} \\
$y=\dfrac{k}{x}$ & triple $x$ & \ans{$y$ becomes one third of what it was} \\
$y=\dfrac{k}{x}$ & cut $x$ in half & \ans{$y$ is doubled} \\
$y=kxz$ & double $x$ \emph{and} triple $z$ & \ans{$y$ becomes $6$ times as big} \\
$y=\dfrac{kx}{z}$ & double $x$ \emph{and} double $z$ & \ans{$y$ does not change at all} \\
\hline
\end{tabularx}
\end{center}
\vspace{-2pt}
Prove one of them instead of trusting it. Start from $y=\dfrac{k}{x}$ and replace $x$ with $3x$:
\[ \text{new } y=\frac{k}{{\color{keyred}\mathbf{3x}}}=\frac{1}{{\color{keyred}\mathbf{3}}}\cdot
   \frac{k}{x}=\frac{1}{{\color{keyred}\mathbf{3}}}\cdot(\text{old } y) \]

\vspace{5pt}
\textbf{Part 3 --- A direct variation in disguise, and a design task.}
\begin{enumerate}[label=(\alph*), leftmargin=2.1em, itemsep=5pt]
  \item Ravi says: ``$y=\dfrac{6}{x}$ is an \emph{inverse} variation between $y$ and $x$ --- but it is
        a \emph{direct} variation between $y$ and $\frac1x$.'' Test him.
        \begin{center}
        \renewcommand{\arraystretch}{1.3}
        \begin{tabular}{|l|c|c|c|c|}
        \hline
        \rowcolor{forestbg}
        $x$ & $1$ & $2$ & $3$ & $6$ \\ \hline
        $\frac1x$ & $1$ & $0.5$ & \ans{$\frac13$} & \ans{$\frac16$} \\ \hline
        $y=\frac6x$ & $6$ & $3$ & \ans{$2$} & \ans{$1$} \\ \hline
        $y\div\frac1x$ & \ans{$6$} & \ans{$6$} & \ans{$6$} & \ans{$6$} \\ \hline
        \end{tabular}
        \end{center}
        \vspace{-2pt}
        Is Ravi right? \ans{yes} \; What is the constant? \ans{$6$} \; So the graph of $y$
        against $\frac1x$ would be a \ans{line through the origin}.
  \item Write an \emph{inverse} variation whose graph passes through $(4,5)$: \ans{$y=\dfrac{20}{x}$}
        \\
        Write a \emph{direct} variation whose graph passes through $(4,5)$: \ans{$y=1.25x$}
  \item Now try to write a direct variation whose graph passes through $(0,3)$. \ans{impossible}
        \\
        Explain what goes wrong, in one sentence. \ansline{$y=kx$ always gives $y=0$ when $x=0$, so
        every direct variation passes through the origin and none of them can also pass through
        $(0,3)$}
  \item Every inverse variation $y=\frac kx$ is the Lesson 4.4 parent $\frac1x$ after one
        transformation. Which one? \ans{a vertical stretch by a factor of $k$} \; And its asymptotes
        are always $x=$ \ans{$0$} and $y=$ \ans{$0$}, no matter what $k$ is --- explain why $k$ cannot
        move them. \ansline{a vertical stretch only rescales outputs: $x=0$ is still the one input with
        no output, and $k$ times something approaching $0$ still approaches $0$}
\end{enumerate}
\end{tcolorbox}

\begin{teachernote}
\textbf{Tier R, Part 1 is the whole skill}, and the two rows are there on purpose: a student who fills
in only the row matching their guess has not tested anything. Table (c) is the one to watch --- $y$
decreases, which draws ``inverse'' from a lot of groups until the product row ($20, 34, 42, 44$) is
actually written down. Note that the products \emph{almost} stabilize, which is exactly why the rule
is ``every row.''

\vspace{3pt}
\textbf{Tier R, Part 3 mirrors the Hook} and is the tier's argument. Both twins climb by $5$ for every
$2$; nothing about the \emph{pattern} distinguishes them. Twin (b) is $y=2.5x+3$ --- a head start of
$3$ --- and the quotient row is the only thing that catches it. If a group finishes early, ask them to
predict $y$ at $x=0$ for both tables before you say anything.

\vspace{3pt}
\textbf{Tier A, Part 2 is the A2.F.2f item.} Push for the sentences, not the numbers: $(6,6)$ is the
square, and the right-hand tail says a rectangle of fixed area can be made arbitrarily short but never
flat. Part (e) is the 4.6 bridge --- $b=-4$ is a perfectly good solution of $bh=36$ and is rejected by
the \emph{situation}, not by the algebra. Do not let ``extraneous'' get used here.

\vspace{3pt}
\textbf{Tier A, Part 3 is the signature error of the lesson.} The student's arithmetic is flawless;
they simply stopped after one row. Expect a group to defend them (``but $k$ \emph{is} $3$ for that
pair''), which is the right conversation: every pair has a quotient and a product, so a single pair
carries no information about the type at all.

\vspace{3pt}
\textbf{Tier E, Part 2 is the most valuable page here} if a group is flying --- it is the whole lesson
without a single number, and the little proof at the bottom makes ``triple $x$, get a third of $y$''
something they can show rather than recite. \textbf{Part 3(a)} is the deepest idea in the lesson: an
inverse variation \emph{is} a direct variation with respect to $\frac1x$, which is why $y=\frac kx$ is
the parent $\frac1x$ scaled by $k$. \textbf{Part 3(c)} usually needs one prompt --- ``what does $y=kx$
give you when $x=0$?'' --- and then it lands.

\vspace{3pt}
\textbf{Joint and combined variation are enrichment} (A2.F.1d names direct and inverse only). Part 1(b)
is worth the time even with a mixed group: they derive $V=\pi r^2h$ rather than being handed it.
\end{teachernote}

\end{document}
